\documentclass[10pt]{beamer}

\usepackage{verbatim}
\usepackage[latin1]{inputenc}
\usepackage{amssymb,amsmath,amsfonts,url}
\usepackage{graphicx,color}
\usepackage{pgf,pgfarrows,pgfnodes,pgfautomata,pgfheaps,pgfshade}

\usepackage[all]{xy}
\usepackage{amstext}
\usepackage{xspace}
\usepackage{mdwtab}
\usepackage[french]{babel}
\usepackage[utf8]{inputenc}


\mode<article> {
  \usepackage{fullpage}
  \usepackage{hyperref}
  \setjobnamebeamerversion{advancedcourse.beamer}
}

\mode<presentation> {
    \usepackage{beamerthemeBoadilla}

  \beamertemplatetransparentcovereddynamic
  \beamertemplateballitem
}


\newcommand{\alertonce}[1]{\addtocounter{beamerpauses}{-1}\alert<+>{#1}}
\newcommand{\alerttwice}[1]{\addtocounter{beamerpauses}{-1}\alert<+-+(1)>{#1}}

\newcommand{\subsubsubsection}[1]{\noindent\textbf{#1}}

\newcommand{\cfbox}[1]{\begin{center}\fbox{#1}\end{center}}
\newcommand{\cbox}[1]{\begin{center} #1 \end{center}}

%%%%%%%%%%%%%%%%%%%%%%%%%%%%%%%%%%%%%%%%%%%%%%%%%%%%%%%%%%%%%%%%%%%%%%%%%%%%%%%%%%%%%%%%%

\graphicspath{{Graphics/}}
%\input def_0.tex



\title[IN100 Cours 8]{\textbf{Fondements informatiques I}\\
Cours 8: Modules et fichiers}

\author[]{Sorina Ionica \texttt{sorina.ionica@uvsq.fr} \\~\\Sandrine Vial  \texttt{sandrine.vial@uvsq.fr} }


\institute{}

\date{}


%\setbeamercolor{normal text}{fg=black}

\begin{document}


\frame{\titlepage}




\section{Les modules}

\frame[containsverbatim]{
\frametitle{Modules}



On veut programmer une application calculatrice. Comment structurer notre programme? 


\begin{itemize}
\item Un fichier contenant tous les fonctions permettant de faire des calculs: racines carr\'ees, puissance, log etc. 
\item Un fichier contenant toutes les fonctions permettant de dessiner l'interface graphique :  dessiner des buttons, des zones de texte etc. 
\item Un programme principale depuis lequel on appelle toutes cettes fonctions 
\end{itemize}

\begin{alertblock}{Solution}
\begin{itemize} 
\item On range les d\'efinitions de fonctions se rapportant \`a une m\^eme application au sein d'un fichier (script) commun baptis\'e {\color{red} module}.
\item Un module est sauvegard\'e sous forme d'un fichier  \texttt{<nom du module>.py}.
\end{itemize}
\end{alertblock}


}


\frame[containsverbatim]{
\frametitle{Modules - Premi\`ere m\'ethode d'importation}


\begin{alertblock}{Utiliser un module}
\begin{itemize}
\item Pour utiliser un module, on se sert de l'instruction \\
\begin{center}
 \texttt{import <nom du module>}
 \end{center}
\item L'ex\'ecution de cette instruction consiste \`a ex\'ecuter le script d\'efinissant le module (ce script peut contenir des instructions autres que des d\'efinitions de fonctions).
\item Pour importer un module, Python a besoin de conna\^itre le chemin qui permet d'acc\'eder au fichier correspondant. Ce chemin doit appara\^itre dans la liste des chemins possibles stock\'es dans la variable \texttt{path} du module \texttt{sys}.
\end{itemize}
\end{alertblock}



\begin{exampleblock}{Exemple}
\small{
\begin{verbatim}
>>> import random
>>> random.randint(0,10)
9
\end{verbatim}
}
\end{exampleblock}

%\begin{itemize}
%\item L'instruction \texttt{import} vous permet d'importer toutes les fonctions du module \texttt{random}
%\item Ensuite, nous utilisons la fonction (ou m\'ethode) \texttt{randint(a,b)} du module random; attention cette fonction renvoie un nombre entier al\'eatoirement entre a inclus et b inclus.
%\end{itemize}

}

\frame[containsverbatim]{
\frametitle{Modules - importation d'une fonction}




\medskip
\begin{alertblock}{}
Dans certains cas, on a besoin de disposer d'une seule  fonction du module. 
\begin{verbatim}
from [module] import [fonction]
\end{verbatim}
\end{alertblock}


\vfill
\begin{exampleblock}{Exemple}
\begin{verbatim}
from math import sqrt
racine = sqrt(49)
\end{verbatim}
\end{exampleblock}


}


\frame[containsverbatim]{
\frametitle{Modules - importation de toutes les  fonctions}




\medskip
\begin{alertblock}{}
 Pour disposer de toutes les fonctions d'un module:
 \begin{center}
\begin{verbatim}
from [module] import *
\end{verbatim}
\end{center}
\end{alertblock}

\begin{exampleblock}{Exemple}
\begin{verbatim}
from math import *
racine = sqrt(49)
angle = pi/6
print(sin(angle))
\end{verbatim}

\end{exampleblock}

\vspace{0.5 cm}

Cette m\'ethode est \`a \'eviter, car on peut importer sans le savoir des noms de fonctions qu'on a d\'ej\`a  d\'efini par ailleurs. 

}


\frame[containsverbatim]{
\frametitle{Modules courants}

\begin{itemize}
\item \texttt{math} : fonctions et constantes math\'ematiques de base (sin, cos, exp, pi...).
\item \texttt{sys} : passage d'arguments, gestion de l'entr\'ee/sortie standard etc...
\item \texttt{os} : dialogue avec le syst\`eme d'exploitation.
\item \texttt{random} : g\'en\'eration de nombres al\'eatoires.
\item \texttt{time} : permet d'acc\'eder aux fonctions g\'erant le temps.
\item \texttt{calendar} : fonctions de calendrier.
\item \texttt{urllib} : permet de r\'ecup\'erer des donn\'ees sur internet depuis python.
\item \texttt{Tkinter} : interface python avec Tk (permet de cr\'eer des objets graphiques; n\'ecessite d'installer Tk).
\item \texttt{re} : gestion des expressions r\'eguli\`eres.
\item \texttt{pickle} : \'ecriture et lecture de structures python (comme les dictionnaires par exemple).
\item \dots 
\end{itemize}
}

\frame[containsverbatim]{
\frametitle{Obtenir de l'aide sur les modules import\'es}

\begin{itemize}
\item Pour obtenir de l'aide sur un module rien de plus simple, il suffit d'invoquer la commande help : 

\medskip
\begin{exampleblock}{}
\begin{verbatim}
 import random
 help(random)
\end{verbatim}
\end{exampleblock}

\bigskip\item Il est aussi possible d'invoquer de l'aide sur une fonction particuli\`ere d'un module en la pr\'ecisant de la mani\`ere suivante : 
\medskip
\begin{exampleblock}{}
\begin{verbatim}
help(random.randint)
\end{verbatim}
\end{exampleblock}

\end{itemize}

}




\end{document}

\section{Les fichiers}

\frame[containsverbatim]{
\frametitle{Utilisation de fichiers}

\begin{itemize}
\item Il est important de dissocier les donn\'ees des programmes qui les utilisent en rangeant ces donn\'ees dans des fichiers s\'epar\'es.

\bigskip\item Le module \texttt{os} contient des fonctions qui permettent de localiser les fichiers :
\begin{itemize}
\item \texttt{getcwd()}: Retourne le chemin du r\'epertoire courant
\medskip \item \texttt{chdir(<ch>)}: Change le r\'epertoire courant qui prend la valeur donn\'ee par la cha\^ine de caract\`eres \texttt{<ch>}
\medskip \item  \texttt{path.isfile(<ch>)}: Retourne un bool\'een qui indique s'il existe un fichier dont le chemin est la cha\^ine de caract\`eres \texttt{<ch>}
\medskip \item  \texttt{path.isdir(<ch>)}: Retourne un bool\'een qui indique s'il existe un r\'epertoire dont le chemin est la cha\^ine de caract\`eres \texttt{<ch>}
\end{itemize}
\end{itemize}

\vfill
\begin{exampleblock}{Sous Linux}
\begin{verbatim}
from os import chdir
chdir("/home/mon_user/exercices")
\end{verbatim}
\end{exampleblock}
}

\frame[containsverbatim]{
\frametitle{Quelques exemples}

\begin{exampleblock}{}
\begin{verbatim}
import os

rep_cour = os.getcwd()
print(rep_cour)
print(path.isfile("test.py"))
\end{verbatim}
\end{exampleblock}

\vfill
\begin{exampleblock}{}
\begin{verbatim}
from os import getcwd, path

rep_cour = getcwd() 
print("le nom du repertoire est",rep_cour)
print(path.isfile("test.py"))
\end{verbatim}
\end{exampleblock}


}

\frame[containsverbatim]{
\frametitle{Utilisation de fichiers}


\begin{itemize}
\item Pour utiliser un fichier identifi\'e par le chemin \texttt{ch} dans un programme Python, il faut commencer par l'ouvrir par l'appel de fonction
\begin{center}
 \texttt{open(<ch>, [<mode>] )}
 \end{center}
  qui retourne un objet de type \texttt{file}. 

\bigskip \item Le param\`etre facultatif \texttt{<mode>} indique le mode d'ouverture du fichier :
\begin{itemize}
\medskip\item \texttt{r} : mode lecture (le fichier doit exister pr\'ealablement)
\medskip\item \texttt{w} : mode \'ecriture (si le fichier existe, les donn\'ees sont \'ecras\'ees, sinon le fichier est cr\'e\'e)
\medskip\item \texttt{a} : mode ajout (si le fichier existe, les donn\'ees \'ecrites vont l'\^etre apr\`es celles existantes, sinon le fichier est cr\'e\'e)
\end{itemize}

\bigskip \item Si le mode est omis, le mode par d\'efaut est \texttt{r}.
\end{itemize}
}


\frame[containsverbatim]{
\frametitle{Utilisation de fichiers}

 Un objet de type \texttt{file} est associ\'e \`a des attributs et des m\'ethodes. \\
 
 En voici quelques-unes :

\begin{itemize}
\medskip\item \texttt{read([<n>])} : Retourne la cha\^ine des \texttt{<n>} caract\`eres restants.
\medskip\item \texttt{write(<ch>)} : Ecrit la cha\^ine de caract\`eres \texttt{<ch>}.
\medskip\item \texttt{close()} : Ferme le fichier quand il est fini d'\^etre utilis\'e.
\medskip\item \texttt{seek(<n>)} : Choisit le caract\`ere \texttt{<n>} comme position courante du fichier.
\medskip\item \texttt{tell()}: Retourne le caract\`ere en position courante.
\end{itemize}
}

\frame[containsverbatim]{
\frametitle{Exemple}


Cr\'eez un fichier dans un \'editeur de texte que vous sauverez dans votre r\'epertoire. 

\medskip
\begin{block}{exemple.txt}
\small{
\begin{verbatim}
Ceci est la premiere ligne
Ceci est la deuxieme ligne
Ceci est la troisieme ligne
Ceci est la quatrieme et derniere ligne\end{verbatim}
}
\end{block}

\begin{exampleblock}{}
\small{
\begin{verbatim}
fichier = open('exemple.txt','r')
print(fichier) 


fichier.readlines()

fichier.close()
print(filin)

fichier.readlines()
\end{verbatim}
}
\end{exampleblock}


}


\frame[containsverbatim]{
\frametitle{Lire dans un fichier}


\begin{exampleblock}{Exemple}
\small{
\begin{verbatim}
filin = open('exemple.txt','r')
lignes  = filin.readlines()
for i in lignes:
    print(i)

filin.close()
\end{verbatim}
}
\end{exampleblock}

\vspace{0.5 cm}


\textbf{\`A retenir:} La fonction \texttt{readlines()} renvoie un tableau contenant les cha\^ines de caract\`eres des lignes.

\vspace{0.5 cm}



\begin{block}{}
\small{
\begin{verbatim}
Ceci est la premiere ligne
Ceci est la deuxieme ligne
Ceci est la troisieme ligne
Ceci est la quatrieme et derniere ligne\end{verbatim}
}
\end{block}


}


\frame[containsverbatim]{
\frametitle{Lire dans un fichier}


\begin{exampleblock}{Exemple}
\small{
\begin{verbatim}
filin = open("exemple.txt","r")
filin.read()

filin.close()
\end{verbatim}
}
\end{exampleblock}


\medskip
\begin{block}{\texttt{exemple.txt}}
\small{
\begin{verbatim}
Ceci est la premiere ligne
Ceci est la deuxieme ligne
Ceci est la troisieme ligne
Ceci est la quatrieme et derniere ligne\end{verbatim}
}
\end{block}

\vspace{0.5 cm}
La fonction \texttt{read()} renvoie une cha\^ine de caract\`eres: \\
\small{
\begin{verbatim}
"Ceci est la premiere ligne\nCeci est la deuxieme ligne\n
Ceci est la troisieme ligne\nCeci est la quatrieme et derniere ligne\n"
\end{verbatim}
}
}


\frame[containsverbatim]{
\frametitle{Parcourir un fichier}



\begin{exampleblock}{Exemple}
\small{
\begin{verbatim}
filin = open("exemple.txt","r")
filin.tell() #0L

filin.readline() #'Ceci est la premiere ligne\n'
filin.tell() #27L
filin.seek(0)
filin.tell() #0L
filin.readline() #'Ceci est la premiere ligne\n'
filin.close()
\end{verbatim}
}
\end{exampleblock}
\vspace{0.5 cm}

La fonction \texttt{readline} lit la ligne \`a partir du caract\`ere courant et renvoie une cha\^ine de caract\`eres. 

}


\frame[containsverbatim]{
\frametitle{Exemple}

\begin{exampleblock}{}
\small{
\begin{verbatim}
animaux = ['girafe', 'hippopotame', 'singe', 'dahu' , 'ornithorynque']
filout = open('exemple2.txt','w')
for i in animaux:
        filout.write(i)

filout.close()

$ more exemple2.txt
girafehippopotamesingedahuornithorynque
$
\end{verbatim}
}
\end{exampleblock}


}


